\documentclass[12pt,a4paper]{article}
\usepackage[UTF8]{ctex}          % 中文支持
\usepackage{verbatim}
\usepackage{fancyhdr}
\usepackage{graphicx}
\usepackage{amsfonts}
\usepackage{booktabs}
\usepackage{amsmath,amssymb}     % 数学符号
\usepackage{listings}            % 代码排版
\usepackage{xcolor}              % 颜色
\usepackage{geometry}            % 页面边距
\usepackage{hyperref}            % 超链接
\geometry{left=1.5cm,right=1.5cm,top=2.5cm,bottom=2.5cm}

% MATLAB 代码样式
\lstset{
    language=Matlab,
    basicstyle=\ttfamily\small,
    keywordstyle=\color{blue},
    commentstyle=\color{gray},
    stringstyle=\color{red},
    numbers=left,
    numberstyle=\tiny\color{gray},
    stepnumber=1,
    numbersep=5pt,
    backgroundcolor=\color{white},
    showspaces=false,
    showstringspaces=false,
    frame=single,
    tabsize=4,
    captionpos=b,
    breaklines=true,
    breakatwhitespace=false,
    escapeinside={\%*}{*)},
    morekeywords={u,v,w,RichardsonExtrapolation,romberg,sinx_over_x,stable_b,exp_m1_over_t,print_romberg}
}

\title{第八周上机作业}
\author{PB23000270 李佩优}
\date{\today}

\begin{document}

\maketitle

\section{基于龙贝格外推的数值积分}

\subsection{程序概述}

程序 \texttt{hw8.m} 实现了龙贝格（Romberg）积分算法，用于计算以下三个在奇点或无穷区间上需要特殊处理的定积分：
\[
\begin{aligned}
I_a &= \int_0^1 \frac{\sin x}{x}\,dx,\\[4pt]
I_b &= \int_{-1}^1 \frac{\cos x - e^x}{\sin x}\,dx,\\[4pt]
I_c &= \int_1^\infty \frac{1}{x\,e^{x}}\,dx \quad \left(\text{通过变换 }x=\frac{1}{t}\text{ 化为 }\int_0^1 \frac{e^{-1/t}}{t}\,dt\right).
\end{aligned}
\]
程序首先对每个积分的被积函数在奇点处进行特殊处理（极限延拓或恒等变形），然后调用统一的龙贝格函数生成外推阵列，最终输出阵列及对角线近似值。阵列行数设为 \(n=7\)。

\subsection{算法原理}

\subsubsection{龙贝格积分方法}
龙贝格积分建立在复化梯形公式的基础上。将区间 \([a,b]\) 逐次二等分，记步长 \(h_k = (b-a)/2^{k-1}\)，得到梯形序列
\[
T_{1,1} = \frac{h_1}{2}\bigl[f(a)+f(b)\bigr],\qquad 
T_{k,1} = \frac{1}{2}T_{k-1,1} + h_k\sum_{i=1}^{2^{k-2}} f\bigl(a+(2i-1)h_k\bigr),\ k\ge 2.
\]
若被积函数充分光滑，梯形值满足渐近展开
\[
I - T_{k,1} = c_2 h_k^2 + c_4 h_k^4 + c_6 h_k^6 + \cdots,
\]
其中只含步长的偶次幂。利用该展开进行理查森（Richardson）外推，可逐列消去低阶误差项：
\[
T_{k,j} = \frac{4^{j-1} T_{k,j-1} - T_{k-1,j-1}}{4^{j-1} - 1},\qquad j=2,3,\dots,k.
\]
由此构造出一个下三角矩阵（龙贝格阵列），其对角线元素 \(T_{n,n}\) 通常具有极高精度，且收敛速度远快于原始梯形序列。

\subsubsection{奇点处理与变量变换}
\begin{itemize}
    \item \textbf{积分 \(I_a\)}：被积函数 \(\frac{\sin x}{x}\) 在 \(x=0\) 处可去奇点的极限为 \(1\)，程序通过向量化掩码令该点直接返回 \(1\)。
    \item \textbf{积分 \(I_b\)}：函数在 \(x=0\) 处需谨慎处理，避免分子分母同时趋于零带来的舍入误差。利用恒等式
    \[
    \cos x - e^x = -2\sin^2\frac{x}{2} - (e^x-1),\qquad \sin x = 2\sin\frac{x}{2}\cos\frac{x}{2},
    \]
    并采用 \texttt{expm1} 以提高小变量处的精度。\(x=0\) 处的极限值为 \(-1\)。
    \item \textbf{积分 \(I_c\)}：原积分区间为 \([1,\infty)\)，作变量代换 \(x = 1/t\)，得到
    \[
    \int_1^\infty \frac{1}{x e^x}\,dx = \int_0^1 \frac{e^{-1/t}}{t}\,dt.
    \]
    被积函数在 \(t=0\) 的极限为 \(0\)，同样通过掩码处理。
\end{itemize}

\subsection{程序结构说明}

程序的主流程包含三个被积函数的定义，调用 \texttt{romberg} 函数计算 7 行龙贝格阵列，最后用 \texttt{print\_romberg} 打印下三角阵列与对角线结果。核心子函数如下：

\begin{lstlisting}[caption={主程序框架及辅助函数}, label={lst:hw8main}]
clear; clc; format long e;

f_a = @(x) sinx_over_x(x);
f_b = @(x) stable_b(x);
f_c = @(t) exp_m1_over_t(t);

n_rows = 7;
R_a = romberg(f_a, 0, 1, n_rows);
print_romberg(R_a);

R_b = romberg(f_b, -1, 1, n_rows);
print_romberg(R_b);

R_c = romberg(f_c, 0, 1, n_rows);
print_romberg(R_c);

function y = sinx_over_x(x)
    y = ones(size(x));
    mask = (x ~= 0);
    y(mask) = sin(x(mask)) ./ x(mask);
end

function y = stable_b(x)
    y = zeros(size(x));
    mask0 = (x == 0);
    y(mask0) = -1;
    x_non = x(~mask0);
    s = sin(x_non/2);
    c = cos(x_non/2);
    y(~mask0) = (-2*s.^2 - expm1(x_non)) ./ (2*s.*c);
end

function y = exp_m1_over_t(t)
    y = zeros(size(t));
    mask0 = (t == 0);
    y(mask0) = 0;
    t_non = t(~mask0);
    y(~mask0) = exp(-1./t_non) ./ t_non;
end
\end{lstlisting}

龙贝格外推核心算法 \texttt{romberg} 的实现如下：

\begin{lstlisting}[caption={龙贝格函数}, label={lst:romberg}]
function R = romberg(f, a, b, n)
    R = zeros(n, n);
    h = b - a;
    R(1,1) = (h/2) * (f(a) + f(b));
    for i = 2:n
        h = h / 2;
        x_new = a + (2*(1:2^(i-2)) - 1) * h;
        sum_f = sum(f(x_new));
        R(i,1) = 0.5 * R(i-1,1) + h * sum_f;
        for j = 2:i
            R(i,j) = (4^(j-1) * R(i,j-1) - R(i-1,j-1)) / (4^(j-1) - 1);
        end
    end
end
\end{lstlisting}

\subsection{数值结果与分析}

程序在北太天元上运行，输出结果如下：

\begin{verbatim}
============= 积分 a: ∫_0^1 sin(x)/x dx =============
龙贝格阵列 (行数 7):
0.9207355 
0.9397933 0.9461459 
0.9445135 0.9460869 0.9460830 
0.9456909 0.9460833 0.9460831 0.9460831 
0.9459850 0.9460831 0.9460831 0.9460831 0.9460831 
0.9460586 0.9460831 0.9460831 0.9460831 0.9460831 0.9460831 
0.9460769 0.9460831 0.9460831 0.9460831 0.9460831 0.9460831 0.9460831 
积分近似值 (对角线):  0.9460831


============= 积分 b: ∫_{-1}^1 (cos x - e^x)/sin x dx =============
龙贝格阵列 (行数 7):
-2.7932067 
-2.3966033 -2.2644022 
-2.2852177 -2.2480892 -2.2470016 
-2.2563259 -2.2466953 -2.2466023 -2.2465960 
-2.2490303 -2.2465984 -2.2465919 -2.2465918 -2.2465917 
-2.2472017 -2.2465921 -2.2465917 -2.2465917 -2.2465917 -2.2465917 
-2.2467442 -2.2465917 -2.2465917 -2.2465917 -2.2465917 -2.2465917 -2.2465917 
积分近似值 (对角线):  -2.2465917


============= 积分 c: ∫_1^∞ (x e^x)^{-1} dx  (变换后 ∫_0^1 e^{-1/t}/t dt) =============
龙贝格阵列 (行数 7):
0.1839397 
0.2273051 0.2417603 
0.2198339 0.2173435 0.2157157 
0.2193510 0.2191900 0.2193131 0.2193702 
0.2193836 0.2193945 0.2194081 0.2194096 0.2194097 
0.2193839 0.2193840 0.2193834 0.2193830 0.2193829 0.2193828 
0.2193839 0.2193839 0.2193839 0.2193839 0.2193839 0.2193839 0.2193839 
积分近似值 (对角线):  0.2193839
\end{verbatim}

\textbf{分析}：
\begin{itemize}
    \item \textbf{积分 \(I_a\)}：该积分即为正弦积分 \(\mathrm{Si}(1)\)，精确值约为 \(0.946083070367183\cdots\)。龙贝格阵列从第一行的 \(0.9207\) 到对角线的 \(0.94608307036718\)，仅用 7 行就已达到 \(10^{-14}\) 量级的精度，显示出极快的收敛速度。从阵列中可以看到，随着外推阶数增加，同一列中的数值迅速趋于稳定。
    \item \textbf{积分 \(I_b\)}：被积函数在 \(x=0\) 附近变化剧烈，但经过稳定处理后，龙贝格阵列的对角线收敛至 \(-2.24659172072862\)，前 7 位有效数字已经稳定。与初始梯形值 \(-2.793\) 相比，精度大幅提升。阵列第二行以后的对角线元素几乎不再变化，表明外推已基本消除误差主导项。
    \item \textbf{积分 \(I_c\)}：该积分本质上是指数积分 \(E_1(1)\)，其标准值约为 \(0.219383934395\cdots\)。对角线近似值 \(0.21938394142680\) 与真值相差约 \(7\times10^{-9}\)，精度同样令人满意。注意到阵列的非对角线元素在外推过程中略有波动（如第 6 行第 3、4 列），这是梯形序列尚未完全进入渐近区域时的正常现象；对角线则继续保持单调稳定。
\end{itemize}

三个积分的结果均表明，仅需 7 次区间剖分（对应 7 行龙贝格阵列），外推即可将精度提升至接近机器精度。龙贝格方法对于仅有可去奇点或通过解析延拓光滑化的被积函数，仍然能保持高阶收敛特性。

\subsection{总结}

本实验实现并验证了龙贝格外推积分算法在处理具有可去奇点、无穷区间或数值敏感型被积函数时的有效性。通过合理处理边界奇点（取极限值或恒等变形），并结合理查森外推消除梯形误差展开中的低阶项，可以在较少计算量下获得高精度数值积分。实验同时表明，对于充分光滑的被积函数，龙贝格阵列的对角线元素收敛极快，是实际数值积分中一种可靠且高效的方法。

\newpage
\section{计算机代码}

\small
\hrule
\verbatiminput{hw8.m}
\hrule

\end{document}